\documentclass[12pt,a4paper]{article}
\usepackage[utf8]{inputenc}
\usepackage[vietnamese]{babel}
\usepackage{amsmath,amssymb}
\usepackage[margin=2cm]{geometry}
\usepackage{enumitem}
\usepackage{fontspec}
\setmainfont{Times New Roman}

\title{\textbf{BÀI TẬP VỀ NHÀ: BÀI 6 -- NHIỆT NÓNG CHẢY RIÊNG}\\
\large Vật Lý 12 -- Chương trình GEMS}
\author{Giáo viên: Kha Khung Hiệp}
\date{}

\begin{document}
\maketitle
\textit{Thời gian làm bài: 45 phút}

\section*{PHẦN I: CÂU HỎI TRẮC NGHIỆM NHIỀU LỰA CHỌN (18 CÂU)}

\textbf{Câu 1:} Đơn vị đo của nhiệt nóng chảy riêng của vật rắn trong hệ SI là:
\begin{enumerate}[label=\Alph*.]
  \item $\text{J/kg} \cdot \text{K}$.
  \item $\text{J/kg}$.
  \item $\text{J}$.
  \item $\text{W/kg}$.
\end{enumerate}

\textbf{Câu 2:} Nhiệt lượng cần cung cấp để làm nóng chảy hoàn toàn một khối lượng chất rắn kết tinh ở nhiệt độ nóng chảy được xác định theo hệ thức:
\begin{enumerate}[label=\Alph*.]
  \item $Q = \lambda \cdot m$.
  \item $Q = m \cdot c \cdot \Delta t$.
  \item $Q = \dfrac{\lambda}{m}$.
  \item $Q = \lambda \cdot m \cdot \Delta t$.
\end{enumerate}

\textbf{Câu 3:} Nhiệt nóng chảy riêng của đồng là $\lambda = 1{,}8 \cdot 10^5\text{ J/kg}$. Phát biểu nào sau đây là đúng về ý nghĩa vật lý của con số này?
\begin{enumerate}[label=\Alph*.]
  \item Khối đồng sẽ tỏa ra nhiệt lượng $1{,}8 \cdot 10^5\text{ J}$ khi đông đặc hoàn toàn.
  \item Mỗi kilôgam đồng cần thu vào nhiệt lượng $1{,}8 \cdot 10^5\text{ J}$ để hóa lỏng hoàn toàn ở nhiệt độ nóng chảy.
  \item Cần cung cấp công suất $1{,}8 \cdot 10^5\text{ W}$ để nung nóng đồng chảy lỏng.
  \item Mỗi kilôgam đồng tỏa ra nhiệt lượng $1{,}8 \cdot 10^5\text{ J}$ khi nóng chảy ở nhiệt độ nóng chảy.
\end{enumerate}

\textbf{Câu 4:} Để xác định nhiệt nóng chảy riêng của nước đá bằng thực nghiệm, ta không cần dùng đến dụng cụ nào sau đây?
\begin{enumerate}[label=\Alph*.]
  \item Cân điện tử.
  \item Oát kế đo điện năng và thời gian.
  \item Lực kế.
  \item Nhiệt kế điện tử.
\end{enumerate}

\textbf{Câu 5:} Điều nào sau đây là sai khi nói về nhiệt nóng chảy của vật rắn?
\begin{enumerate}[label=\Alph*.]
  \item Nhiệt nóng chảy của vật rắn là nhiệt lượng cung cấp cho vật rắn trong suốt quá trình nóng chảy.
  \item Đơn vị của nhiệt nóng chảy là Jun ($\text{J}$).
  \item Các chất khác nhau có khối lượng bằng nhau thì nhiệt nóng chảy ở cùng nhiệt độ nóng chảy là như nhau.
  \item Nhiệt nóng chảy được tính bằng công thức $Q = \lambda \cdot m$.
\end{enumerate}

\textbf{Câu 6:} Nhiệt nóng chảy riêng của vật rắn kết tinh phụ thuộc vào:
\begin{enumerate}[label=\Alph*.]
  \item bản chất của vật rắn và áp suất ngoài.
  \item hình dạng và kích thước của vật rắn.
  \item nhiệt độ ban đầu của vật rắn.
  \item tốc độ truyền nhiệt của lò nung.
\end{enumerate}

\textbf{Câu 7:} Ở áp suất tiêu chuẩn, nước đá có nhiệt độ nóng chảy bằng:
\begin{enumerate}[label=\Alph*.]
  \item $0^{\circ}\text{C}$ ($273\text{ K}$).
  \item $100^{\circ}\text{C}$ ($373\text{ K}$).
  \item $0\text{ K}$ ($-273^{\circ}\text{C}$).
  \item $37^{\circ}\text{C}$ ($310\text{ K}$).
\end{enumerate}

\textbf{Câu 8:} Trong đồ thị biểu diễn sự thay đổi nhiệt độ theo thời gian nhận nhiệt của chất rắn kết tinh, giai đoạn chất rắn đang nóng chảy được biểu diễn bởi:
\begin{enumerate}[label=\Alph*.]
  \item đường cong đi lên.
  \item đoạn thẳng nằm ngang song song với trục thời gian.
  \item đoạn thẳng dốc đi xuống.
  \item đường parabol hướng lên.
\end{enumerate}

\textbf{Câu 9:} Nhiệt lượng cần cung cấp để làm nóng chảy hoàn toàn cục nước đá có khối lượng $m = 200\text{ g}$ ở $0^{\circ}\text{C}$ là bao nhiêu? Biết nhiệt nóng chảy riêng của nước đá là $\lambda = 3{,}4 \cdot 10^5\text{ J/kg}$.
\begin{enumerate}[label=\Alph*.]
  \item $6{,}8 \cdot 10^4\text{ J}$.
  \item $6{,}8 \cdot 10^7\text{ J}$.
  \item $1{,}7 \cdot 10^6\text{ J}$.
  \item $6{,}8 \cdot 10^5\text{ J}$.
\end{enumerate}

\textbf{Câu 10:} Trong thí nghiệm đo nhiệt nóng chảy riêng của nước đá, oát kế kỹ thuật số đo trực tiếp đại lượng nào sau đây?
\begin{enumerate}[label=\Alph*.]
  \item Nhiệt nóng chảy riêng của nước đá.
  \item Công suất đun dây điện trở và thời gian dòng điện chạy qua.
  \item Khối lượng của nước đá tan.
  \item Độ tăng nhiệt độ của hỗn hợp.
\end{enumerate}

\textbf{Câu 11:} Hệ thức liên hệ giữa điện năng tiêu thụ từ nguồn điện nung $A = P \cdot t$ và nhiệt lượng nước đá nhận vào trong trường hợp bỏ qua hao phí nhiệt lượng truyền ra môi trường là:
\begin{enumerate}[label=\Alph*.]
  \item $P \cdot t = \lambda \cdot m$.
  \item $P \cdot t = m \cdot c \cdot \Delta t$.
  \item $P \cdot t = \dfrac{\lambda}{m}$.
  \item $\lambda = \dfrac{m}{P \cdot t}$.
\end{enumerate}

\textbf{Câu 12:} Khi tăng áp suất ngoài tác dụng lên mặt thoáng của chất rắn kết tinh, nhiệt độ nóng chảy của nó thay đổi như thế nào?
\begin{enumerate}[label=\Alph*.]
  \item Luôn luôn tăng đối với mọi vật rắn.
  \item Luôn luôn giảm đối với mọi vật rắn.
  \item Tăng đối với vật dãn nở khi nóng chảy và giảm đối với vật co lại khi nóng chảy.
  \item Giảm đối với vật dãn nở khi nóng chảy và tăng đối với vật co lại khi nóng chảy.
\end{enumerate}

\textbf{Câu 13:} Nhiệt lượng cần cung cấp cho $2\text{ kg}$ nước đá ở $-10^{\circ}\text{C}$ để nó hóa lỏng hoàn toàn thành nước ở $0^{\circ}\text{C}$ là bao nhiêu? Cho biết nhiệt dung riêng của nước đá là $c_{đá} = 2100\text{ J/kg.K}$ và nhiệt nóng chảy riêng của nước đá là $\lambda = 3{,}4 \cdot 10^5\text{ J/kg}$.
\begin{enumerate}[label=\Alph*.]
  \item $7{,}22 \cdot 10^5\text{ J}$.
  \item $6{,}80 \cdot 10^5\text{ J}$.
  \item $4{,}20 \cdot 10^4\text{ J}$.
  \item $7{,}22 \cdot 10^4\text{ J}$.
\end{enumerate}

\textbf{Câu 14:} Khi đổ $350\text{ g}$ thiếc nóng chảy ở nhiệt độ $232^{\circ}\text{C}$ vào một nhiệt lượng kế đựng $330\text{ g}$ nước ở $7^{\circ}\text{C}$. Nhiệt dung của nhiệt lượng kế là $Q_{kế} = 100\text{ J/K}$. Sau khi cân bằng nhiệt, nhiệt độ nước trong bình là $32^{\circ}\text{C}$. Biết nhiệt dung riêng của nước là $c_{nước} = 4200\text{ J/kg.K}$, của thiếc rắn là $c_{thiếc} = 230\text{ J/kg.K}$. Nhiệt nóng chảy của thiếc gần giá trị nào nhất sau đây?
\begin{enumerate}[label=\Alph*.]
  \item $0{,}6 \cdot 10^5\text{ J/kg}$.
  \item $0{,}6 \cdot 10^4\text{ J/kg}$.
  \item $3{,}4 \cdot 10^5\text{ J/kg}$.
  \item $1{,}8 \cdot 10^5\text{ J/kg}$.
\end{enumerate}

\textbf{Câu 15:} Để đúc các chi tiết máy bằng thép, lò nung đốt hết $200\text{ kg}$ than đá có năng suất tỏa nhiệt là $q = 3 \cdot 10^7\text{ J/kg}$. Biết hiệu suất của lò là $30\%$ (chỉ $30\%$ nhiệt lượng lò tỏa ra dùng để nung và nóng chảy thép). Thép đưa vào lò có nhiệt độ ban đầu là $20^{\circ}\text{C}$ và nhiệt độ nóng chảy của thép là $1400^{\circ}\text{C}$. Cho biết thép có nhiệt dung riêng rắn là $460\text{ J/kg.K}$, nhiệt nóng chảy riêng là $\lambda = 2{,}7 \cdot 10^5\text{ J/kg}$. Khối lượng mẻ thép bị nấu chảy có giá trị gần nhất nào sau đây?
\begin{enumerate}[label=\Alph*.]
  \item $2{,}0\text{ tấn}$.
  \item $2{,}0\text{ tạ}$.
  \item $4{,}0\text{ tấn}$.
  \item $4{,}0\text{ tạ}$.
\end{enumerate}

\textbf{Câu 16:} Thả một cục nước đá có khối lượng $m = 30\text{ g}$ ở $0^{\circ}\text{C}$ vào một cốc nước chứa $200\text{ g}$ nước ở $20^{\circ}\text{C}$. Bỏ qua nhiệt dung của cốc và sự truyền nhiệt ra môi trường. Cho biết nhiệt dung riêng của nước là $c = 4200\text{ J/kg.K}$, nhiệt nóng chảy của nước đá là $\lambda = 3{,}4 \cdot 10^5\text{ J/kg}$. Nhiệt độ cuối cùng của cốc nước khi đạt cân bằng nhiệt là:
\begin{enumerate}[label=\Alph*.]
  \item $6{,}3^{\circ}\text{C}$.
  \item $10{,}2^{\circ}\text{C}$.
  \item $6{,}8^{\circ}\text{C}$.
  \item $0^{\circ}\text{C}$ (nhưng đá chưa tan hết).
\end{enumerate}

\textbf{Câu 17:} Phát biểu nào sau đây là sai về sự nóng chảy?
\begin{enumerate}[label=\Alph*.]
  \item Chất rắn kết tinh nóng chảy ở một nhiệt độ xác định dưới áp suất cho trước.
  \item Chất rắn vô định hình không có nhiệt độ nóng chảy xác định.
  \item Khi bắt đầu nóng chảy, dù tiếp tục cung cấp nhiệt lượng, nhiệt độ chất rắn vô định hình vẫn không đổi.
  \item Trong suốt quá trình nóng chảy của chất rắn kết tinh, nhiệt độ của hệ không đổi.
\end{enumerate}

\textbf{Câu 18:} Hỗn hợp gồm $m_1$ nước đá và $m_2$ nước lỏng đựng trong một nhiệt lượng kế cách nhiệt đạt cân bằng nhiệt ở $0^{\circ}\text{C}$. Nếu ta cung cấp thêm một nhiệt lượng $Q = 1{,}7 \cdot 10^4\text{ J}$ thì có $50\text{ g}$ nước đá bị tan chảy. Nhiệt nóng chảy riêng của nước đá bằng:
\begin{enumerate}[label=\Alph*.]
  \item $3{,}4 \cdot 10^5\text{ J/kg}$.
  \item $3{,}4 \cdot 10^4\text{ J/kg}$.
  \item $1{,}7 \cdot 10^5\text{ J/kg}$.
  \item $3{,}4 \cdot 10^6\text{ J/kg}$.
\end{enumerate}

\section*{PHẦN II: CÂU HỎI ĐÚNG/SAI (4 CÂU, MỖI CÂU 4 Ý)}

\textbf{Câu 1:} Trong thí nghiệm xác định nhiệt nóng chảy riêng của nước đá bằng sơ đồ nguyên lý đun nóng điện tích hợp Watt kế:
\begin{enumerate}[label=\alph*)]
  \item Nhiệt lượng kế bằng nhựa xốp nhằm ngăn cản sự trao đổi nhiệt giữa hệ bên trong với môi trường bên ngoài.
  \item Công thức xác định nhiệt nóng chảy thực nghiệm của nước đá khi bỏ qua hao phí là $\lambda = \dfrac{P \cdot t}{m}$.
  \item Oát kế kỹ thuật số được nối nối tiếp với mạch đun dây điện trở để trực tiếp đếm số dòng điện dịch chuyển.
  \item Dùng máy khuấy liên tục hỗn hợp nước đá nhằm đảm bảo nhiệt độ của hệ phân bố đều và quá trình tan diễn ra đồng đều.
\end{enumerate}

\textbf{Câu 2:} Một chậu cách nhiệt chứa hỗn hợp gồm $m_{đá}$ nước đá và $m_{nước}$ nước lỏng tổng khối lượng $10\text{ kg}$ ở $0^{\circ}\text{C}$. Chậu đặt trong phòng ấm và nhiệt độ của hỗn hợp được ghi nhận theo thời gian đun tự nhiên. Đồ thị cho thấy: trong 50 phút đầu nhiệt độ không đổi ở $0^{\circ}\text{C}$, 10 phút sau nhiệt độ tăng từ $0^{\circ}\text{C}$ lên $2^{\circ}\text{C}$. Cho biết $c_{nước} = 4200\text{ J/kg.K}$ và $\lambda_{đá} = 3{,}4 \cdot 10^5\text{ J/kg}$. Bỏ qua nhiệt dung chậu:
\begin{enumerate}[label=\alph*)]
  \item Trong 50 phút đầu, hỗn hợp không thu nhiệt lượng nào từ môi trường ngoài.
  \item Nhiệt lượng nhận được từ môi trường trong 10 phút sau (từ phút thứ 50 đến 60) để tăng nhiệt độ của nước là $84000\text{ J}$.
  \item Tốc độ nhận nhiệt trung bình của chậu từ môi trường là $140\text{ J/s}$.
  \item Khối lượng nước đá ban đầu trong chậu gần bằng $1{,}24\text{ kg}$.
\end{enumerate}

\textbf{Câu 3:} Khảo sát quá trình chuyển thể từ rắn sang lỏng của hai nhóm chất rắn kết tinh và chất rắn vô định hình:
\begin{enumerate}[label=\alph*)]
  \item Thủy tinh, nhựa đường, và sáp nến khi nung nóng sẽ dẻo dần rồi chảy lỏng hoàn toàn mà không có nhiệt độ nóng chảy cố định.
  \item Đồng, sắt, nhôm có cấu trúc mạng tinh thể chặt chẽ nên nóng chảy ở các nhiệt độ xác định.
  \item Trong quá trình làm nóng chảy hoàn toàn thanh nhôm ở nhiệt độ nóng chảy của nó, nếu ta tắt nguồn nhiệt thì nhôm sẽ lập tức đông đặc hóa rắn ngay.
  \item Mặc dù nước đá có $\lambda = 3{,}4 \cdot 10^5\text{ J/kg}$ lớn hơn đồng ($\lambda = 1{,}8 \cdot 10^5\text{ J/kg}$), đồng vẫn cần nhiệt độ cao hơn nhiều để nóng chảy do đồng có liên kết ion/kim loại bền hơn liên kết hydro của nước đá.
\end{enumerate}

\textbf{Câu 4:} Để xác định nhiệt nóng chảy của một kim loại X chưa biết, người thợ nung nóng chảy hoàn toàn $500\text{ g}$ kim loại X ở nhiệt độ nóng chảy $327^{\circ}\text{C}$ rồi đổ nhanh vào cốc đồng khối lượng $200\text{ g}$ chứa $400\text{ g}$ nước ở $20^{\circ}\text{C}$ đặt trong nhiệt lượng kế. Khi đạt cân bằng nhiệt, nhiệt độ đo được là $40^{\circ}\text{C}$. Cho biết $c_{nước} = 4180\text{ J/kg.K}$, $c_{đồng} = 380\text{ J/kg.K}$, $c_{X} = 130\text{ J/kg.K}$:
\begin{enumerate}[label=\alph*)]
  \item Nhiệt lượng cốc đồng và nước nhận vào từ kim loại X tỏa ra là $35120\text{ J}$.
  \item Khi kim loại X nguội đi từ trạng thái rắn ở $327^{\circ}\text{C}$ xuống $40^{\circ}\text{C}$ tỏa ra nhiệt lượng là $18655\text{ J}$.
  \item Nhiệt lượng tỏa ra trong quá trình hóa rắn của kim loại X là $16465\text{ J}$.
  \item Nhiệt nóng chảy riêng của kim loại X tính được bằng thực nghiệm xấp xỉ bằng $0{,}33 \cdot 10^5\text{ J/kg}$.
\end{enumerate}

\section*{PHẦN III: CÂU HỎI TRẢ LỜI NGẮN (6 CÂU)}

\textbf{Câu 1:} Tính nhiệt lượng cần cung cấp (theo đơn vị kilôjun - kJ) để làm nóng chảy hoàn toàn $2\text{ kg}$ nước đá đang ở nhiệt độ $0^{\circ}\text{C}$. Cho biết nhiệt nóng chảy riêng của nước đá là $\lambda = 3{,}4 \cdot 10^5\text{ J/kg}$.\\
\textbf{Đáp số:} \dotfill

\textbf{Câu 2:} Tính nhiệt lượng tổng cộng (theo đơn vị kilôjun - kJ) cần cung cấp để chuyển $0{,}5\text{ kg}$ nước đá từ nhiệt độ $-10^{\circ}\text{C}$ thành nước lỏng hoàn toàn ở nhiệt độ $0^{\circ}\text{C}$. Cho biết nhiệt dung riêng của nước đá là $c_{đá} = 2100\text{ J/kg.K}$ và nhiệt nóng chảy riêng của nước đá là $\lambda = 3{,}4 \cdot 10^5\text{ J/kg}$.\\
\textbf{Đáp số:} \dotfill

\textbf{Câu 3:} Một thỏi nhôm có khối lượng $1\text{ kg}$ ở nhiệt độ ban đầu $20^{\circ}\text{C}$. Biết nhôm có nhiệt độ nóng chảy là $658^{\circ}\text{C}$, nhiệt dung riêng của nhôm rắn là $896\text{ J/kg.K}$ và nhiệt nóng chảy riêng của nhôm là $3{,}9 \cdot 10^5\text{ J/kg}$. Tính nhiệt lượng tối thiểu cần cung cấp để làm thỏi nhôm nóng chảy hoàn toàn (theo đơn vị mêgajun - MJ, kết quả làm tròn đến hai chữ số thập phân).\\
\textbf{Đáp số:} \dotfill

\textbf{Câu 4:} Thả một cục nước đá có khối lượng $30\text{ g}$ ở nhiệt độ $0^{\circ}\text{C}$ vào một cốc nước cách nhiệt chứa $200\text{ g}$ nước ở $20^{\circ}\text{C}$. Bỏ qua nhiệt dung của cốc. Cho biết nhiệt dung riêng của nước là $4{,}2\text{ J/g.K}$ và nhiệt nóng chảy riêng của nước đá là $334\text{ J/g}$. Tính nhiệt độ cuối cùng của cốc nước khi đạt trạng thái cân bằng nhiệt (đơn vị: độ C, kết quả làm tròn đến một chữ số thập phân).\\
\textbf{Đáp số:} \dotfill

\textbf{Câu 5:} Để đúc các vật thể trang trí bằng đồng, người nghệ nhân nấu chảy đồng trong lò nung điện. Lò nung nhận mẻ đồng $10\text{ kg}$ ở nhiệt độ ban đầu $24^{\circ}\text{C}$. Đồng nóng chảy ở $1084^{\circ}\text{C}$, nhiệt dung riêng của đồng rắn là $380\text{ J/kg.K}$, nhiệt nóng chảy riêng của đồng là $1{,}8 \cdot 10^5\text{ J/kg}$. Tính nhiệt lượng có ích mà mẻ đồng hấp thụ để hóa lỏng hoàn toàn (đơn vị: MJ, kết quả làm tròn đến hai chữ số thập phân).\\
\textbf{Đáp số:} \dotfill

\textbf{Câu 6:} Thí nghiệm đo nhiệt nóng chảy của chì: đổ $400\text{ g}$ chì nóng chảy ở nhiệt độ $327^{\circ}\text{C}$ vào một nhiệt lượng kế đựng $300\text{ g}$ nước ở $15^{\circ}\text{C}$. Nhiệt dung của bình cách nhiệt là $120\text{ J/K}$. Nhiệt độ cân bằng cuối cùng là $25^{\circ}\text{C}$. Biết nhiệt dung riêng của nước là $4180\text{ J/kg.K}$, của chì rắn là $130\text{ J/kg.K}$. Hãy tính nhiệt nóng chảy riêng của chì thu được từ thí nghiệm này (đơn vị: $10^4\text{ J/kg}$, kết quả làm tròn đến một chữ số thập phân).\\
\textbf{Đáp số:} \dotfill

\vfill
\begin{center}
\textit{--- GEMS Vật Lý 12 -- Bài 6: Nhiệt Nóng Chảy Riêng ---}
\end{center}

\end{document}